\section{Conclusion}
\label{sec:conclusion}

Progressively evaluating a fixed differential-landmark prefix is correct and
can avoid substantial heuristic work, but the additional OPEN cycle is not
free.  Under the locked $4\rightarrow32$ schedule, staging saved a median
3.10\% of pivot evaluations and 3.02\% of distance reads, while its primary
paired time ratio relative to ordinary Lazy \Astar{} was 1.0543 with a 95\%
hierarchical bootstrap interval of $[0.9890,1.1016]$.  The pooled result does
not support a universal speed advantage.

The map regimes answer the more useful question.  A four-pivot prefix helped
on mazes, where it removed about 30\% of reads and both maps favored staging.
It hurt on warehouses, where median read saving was zero and both maps favored
ordinary lazy evaluation.  Thus the relevant condition is whether the prefix
filters enough full evaluations to repay its extra pops and requeues.  That
finding is descriptive and implementation-conditional, but it is supported by
exact work counters, a negative read-saving/time-ratio association, independent
BFS validation, and identical full-landmark expansion traces.

The study's scientific value lies in disciplined scope: it combines familiar
course concepts into a reproducible experiment, states an honest prior-art
boundary, preserves a sealed map-disjoint protocol, and reports a mixed result
rather than tuning it away.  Future work could preregister multiple prefix
sizes or a deployment selector on new maps, port the engine to a lower-level
language, and test weighted or directed graphs.  Those extensions require new
data and are not inferred from this evaluation.
